\begin{center}
    {\large \textbf{Giorno 8} \textbar\ The Party Island \textbar\ 16 Agosto 2024 \textbar\ \textbf{Gili Trawangan}, Lombok}
\end{center}

\noindent\rule{\textwidth}{0.4pt}
\begin{figure}[h!] % Portrait images below
    \centering
    % First portrait image (on the left)
    \begin{minipage}[h]{0.48\textwidth} % Set width equal to half the page
        \centering
        \includegraphics[width=\textwidth]{images/flags/Screenshot GiliT.png}
    \end{minipage}
    \hfill
    % Text
    \begin{minipage}[h]{0.48\textwidth}
        \raggedright
        \textbf{Superficie}: 4.567 km² \\
        \vspace{0.3cm}
        \textbf{Popolazione}: 3,3 milioni \\
        \vspace{0.3cm}
        \textit{L'isola}
    \end{minipage}
\end{figure}

\landscapeLargeMargin{images/day8/PXL_20240816_001529403.MP.jpg}

\landscapeLargeMargin{images/day8/PXL_20240816_031431646.jpg}

\landscapeLargeMargin{images/day8/PXL_20240816_033610883.jpg}

\landscapeLargeMargin{images/day8/PXL_20240816_035137922.jpg}

\landscapeLargeMargin{images/day8/PXL_20240816_035442750.jpg}

\landscapeLargeMargin{images/day8/PXL_20240816_041017921.jpg}

\landscapeLargeMargin{images/day8/PXL_20240816_041847529.MP.jpg}

\landscapeLargeMargin{images/day8/PXL_20240816_041912775.jpg}

\noindent 16 Agosto 2024 \newline

Notte di riposo, seppur breve. Al risveglio siamo spensierati: le valigie sono pronte, il viaggio dovrebbe procedere senza intoppi e ci aspettano tre giorni dedicati al relax in un paradiso tropicale.

Torniamo a fare colazione al ristorante del Novotel, nuovamente a buffet, ma stavolta rimaniamo più composti e senza sfondarci di cibo, dal momento che nell'ultima settimana ci siamo disabituati a mangiare tanto e qui servono i pasti col badile. Ci limitiamo quindi a un paio di pancake con burro d'arachidi, pain au chocolat, brioche varie, torta Red Velvet, pancake di zucca e per Eli anche un po' di noodles.

Terminata la colazione recuperiamo zaini, borse e lo scatolone con i souvenir di terracotta e ci rechiamo alla reception per il check-out. Ad ogni spostamento siamo un po' più carichi. Consegniamo le chiavi senza intoppi e la nostra macchina arriva puntuale, anche se stavolta l'autista parla a stento inglese. Poco male: oggi non abbiamo molto interesse a fare conversazione.

Trascorriamo il viaggio guardando fuori dal finestrino, scattando foto e trascrivendo i ricordi dei giorni precedenti. Io in particolare continuo a sbobinare le registrazioni dei tour nei villaggi Sasak. Attraversiamo Mataram, la capitale di Lombok, molto caotica e piena di motorini, simile a Jakarta per quel poco che possiamo scrutare dalla macchina in movimento. Procediamo poi per zone collinari immerse nella giungla, e sui lati della strada vediamo numerose bancarelle di nativi che vendono cibo, accessori di telefonia e cianfrusaglie varie. Molto probabilmente quei bugigattoli sono usati anche come abitazioni e non osiamo pensare a cosa succeda loro durante un nubifragio.

Appena accediamo a una zona più aperta il guardrail si popola di macachi: questa volta ce ne sono quasi un centinaio ed Elisa continua a scattare foto come una mitragliatrice.

A un certo punto mi dice dal nulla: "Devo farti assaggiare quel frutto che puzza: l'anduja." Era il durian.

In un paio d'ore raggiungiamo il porto di Bangsal, da cui parte il traghetto per le isole Gili. Il nostro autista ci scarica nel parcheggio del porto nelle mani di un venditore di biglietti, che ci accompagna alla bancarella di un compare. Ci mostrano il listino dei prezzi e ci comunicano che il ferry turistico parte ogni volta che vengono raggiunti quaranta passeggeri e che al momento ve ne sono solo cinque prenotati. In alternativa possiamo prendere il ferry privato, a un prezzo molto più alto.

Considerando che avevo già deciso di non comprare da loro i biglietti non appena eravamo stati approcciati, gli comunico che dobbiamo aspettare degli amici immaginari e che poi saremmo tornati. Esploriamo il porto e in meno di un minuto troviamo l'imbarco ufficiale, dove il biglietto standard costa meno — di poco, ma in ogni caso non siamo qui a farci fregare — e dove vi sono già oltre trenta passeggeri pronti a partire.

Acquistati i biglietti ci rechiamo sulla banchina ad attendere l'imbarco e non facciamo in tempo a sederci che ci chiamano al megafono per salire a bordo. Ci mettiamo scalzi ed entriamo in mare per raggiungere la barca. Non un'impresa semplice con zaini e borsoni, ma con un po' di equilibrio ci riusciamo senza problemi.

Il barcone ondeggia instabile, inclinato costantemente a destra, tanto che alcuni ragazzi si spostano sull'altro lato per bilanciarlo. Sono soprattutto gli indonesiani a prendere l'iniziativa, stipati a poppa insieme a noi, mentre i turisti si trovano a prua. Il viaggio è breve, circa mezz'ora, durante la quale passiamo accanto a Gili Air e Gili Meno. Il caldo è soffocante e per combatterlo apriamo i finestrini, lasciando entrare un po' d'aria. Arriviamo sani e salvi a Gili Trawangan.

Durante il tragitto un motoscafo completamente rosa ci taglia la strada a tutta velocità e già inizio a temere per la tipologia di divertimento che troveremo sull'isola.

Attracchiamo e scendiamo a terra, carichi come muli e con le scarpe in mano. Dal porto cerchiamo di capire come raggiungere l'hotel. Sull'isola non circolano mezzi a motore e per arrivarci a piedi impiegheremmo mezz'ora, ma con tutto il carico che abbiamo proprio non se ne parla. Lascio Elisa con le borse a fare la guardia e vado in avanscoperta nella caotica via principale. Sento già parlare troppo italiano per i miei gusti.

Un ragazzo indonesiano mi approccia subito proponendomi di noleggiare delle biciclette, il mezzo preferito per spostarsi sull'isola. Coi bagagli la vedo complicata, quindi mi propone il secondo mezzo preferito: il calesse trainato dai cavalli. È gentile e ne ferma uno che mi segue fino da Elisa. Carichiamo i bagagli seguendo le precise istruzioni del cocchiere per bilanciare il peso e partiamo. Il cavallo corre come un matto, ma pare che questo sia uno standard qui: schiva pedoni, biciclette e altri cavalli e per fortuna dopo un po' si immette in una stradina secondaria meno affollata, diretta verso il centro dell'isola. Superiamo B\&B, guesthouse, negozietti, bar e ristoranti fino a giungere in una zona meno popolata con terreni liberi e campi incolti. Mi meraviglio inizialmente che il nostro hotel sia in una zona del genere, ma non appena raggiungiamo il Villas Edenia — pronunciato "Edenìa" — mi ricredo subito.

Il posto è meraviglioso: immerso nel verde, regala subito un'atmosfera di relax. Mentre dei ragazzi ci aiutano con i bagagli, il giovane responsabile della reception ci accoglie con un sorriso mentre un suo collega ci consegna un succo d'arancia di benvenuto e degli asciugamani freschi e profumati per rinfrescarci. Sono sempre molto educati e interessati nei convenevoli, e quando scopre che siamo stati sulla Gili Asahan ci comunica che loro sono "brothers" con un resort sull'isola. Mi sembra si tratti proprio del Pearl Beach Resort dove siamo stati noi, ma non ci metto la mano sul fuoco.

Qui al Villas Edenia avevamo una prenotazione per quattro persone. Avremmo dovuto riunirci con Giulia e Ugo, i due amici con i quali avevamo trascorso qualche giorno in Messico, ma per motivi di salute insorti la settimana prima di partire hanno dovuto dare forfait, senza riuscire a cancellare la prenotazione in tempo. Spiego la situazione al check-in, ma a quanto pare sono già aggiornati e non ci sono problemi. La struttura ospita al suo interno un ristorante italiano, "L'Osteria dell'Isola", e ci comunicano che è possibile ordinare tramite un menù digitale e WhatsApp. La notizia bomba è che avendo prenotato per quattro abbiamo la doppia colazione inclusa.

Terminati i convenevoli ci portano alla nostra stanza — anzi, alla nostra villa — attraversando un sentiero circondato da banani e palme da cocco. Apriamo la porta e veniamo travolti dall'effetto "wow". Il posto toglie il fiato. Dal momento che io mi annoio facilmente e che avrei descritto il tutto come "doppia villa con piscina", lascio spazio a Elisa:

\begin{center}
    % Decorative line
    \noindent\rule{0.6\textwidth}{0.4pt}
\end{center}

\textit{"Ci sono due villette private, molto spaziose, dal tetto di paglia, pilastri in bambù e legno sbiancato, vetrate scorrevoli, piscina privata con idromassaggio} (nota di Marco: non era idromassaggio, ma il getto del ricircolo dell'acqua), \textit{patio esterno con divano, tavolino, frigo-bar e vari puff per prendere il sole. In dotazione abbiamo un vassoio in bambù per cibi e bevande da utilizzare direttamente dentro la piscina. Il pavimento in legno è circondato da piante esotiche che creano un'atmosfera unica: banani, palme, bouganville fucsia e corallo, frangipane dai fiori bianchi e gialli con il loro inconfondibile profumo tropicale. Spicca una strelitzia in fiore, con la sua forma che ricorda un uccello esotico, tanto da essere chiamata "uccello del paradiso" per i suoi colori vivaci, blu cangiante e arancione. Completano il quadro gelsomino bianco profumato, piante di avocado, gigantesche orecchie d'elefante (Alocasia macrorrhiza) e l'imponente strelitzia gigante (Alocasia cucullata). I bagni sono all'esterno con la solita scala di bambù per gli asciugamani, la doccia in pietra e i sanitari sotto una tettoia di paglia."}

\begin{center}
    % Decorative line
    \noindent\rule{0.6\textwidth}{0.4pt}
\end{center}

Che altro? Happy hour tutto il pomeriggio con cocktail a metà prezzo, e ovviamente il rifornimento di birra e bibite nel frigo-bar. Leggiamo online che nei supermercati sparsi per l'isola il costo dei prodotti da frigo è lo stesso del frigo-bar della villa, quindi birra a volontà.

Dopo il lungo itinerario decidiamo di trascorrere il pomeriggio tra bagni, sole e relax nella quiete del posto.

Enrico e Cristina stanno arrivando da Bali ma hanno avuto dei ritardi all'imbarco del traghetto, così ce la prendiamo ancora più comoda e ordiniamo dei nachos per merenda. Dopo aver inviato l'ordine mi viene un dubbio atroce e corro alla reception ad assicurarmi che il conto del servizio in camera non venga addebitato ai nostri amici rimasti a casa: ci manca solo che, dopo aver rinunciato al viaggio, si vedano arrivare anche i nostri conti da pagare.

Quando Enrico e Cristina arrivano proseguiamo l'aperitivo tutti insieme. Passando dalla reception per avvisare del loro arrivo conosco Zidane, il nuovo ragazzo alla reception, anche lui molto gentile e alla mano, che da quel momento mi chiama per nome e mi saluta ogni volta che passo. Una volta riuniti ci raccontiamo a vicenda del viaggio finora. Dai loro racconti mi viene il timore che Bali sia una meta troppo commerciale, ma sono certo che troveremo il modo di viverla davvero, come facciamo sempre io ed Elisa.

Dopo quattro chiacchiere e altrettante birre ci salutiamo per cambiarci e prepararci a uscire. Quando usciamo dalla villa l'area circostante è tranquilla e abbastanza frequentata, un punto di passaggio obbligato tra le due coste. Siamo tutti affamati e non conoscendo ancora bene l'isola decidiamo di fermarci in un posto lungo la strada che non sembra troppo turistico e pare avere piatti locali tipici.

Il Gili Frame Resto and Bar si trova in mezzo alla vegetazione, in una struttura aperta di bambù con luminarie colorate e un'atmosfera vagamente reggae. Al tavolo troviamo bustine di Autan da spalmare sulla pelle contro le zanzare. Il nostro cameriere ci scambia per spagnoli — o più probabilmente è la lingua mediterranea che conosce meglio — e ci prende le ordinazioni a suon di "gracias" e "bueno". Io provo finalmente il mio primo Nasi Goreng, Mie Goreng per Elisa e Curry Soup per Enrico e Cristina. Il mio piatto vince: è buono, abbondante e gli spiedini di pollo Satay con la salsa di arachidi valgono da soli il prezzo — molto ragionevole — del piatto.

Terminiamo la cena e proseguiamo la passeggiata fino ad arrivare alla costa est dell'isola, dalla quale intravediamo in lontananza le luci di Gili Meno. La stradina principale è molto animata, con ristoranti, locali sulla spiaggia con musica live, pub e discoteche. Percorriamo tutta la costa fino alle zone più tranquille, per adulti e famiglie, dove ci si può accomodare a sorseggiare un drink sulla spiaggia o a bordo di una delle piscine dei resort. Il posto ha molti aspetti in comune con Palma di Maiorca, Tel Aviv o Rio de Janeiro, ma possiede un'atmosfera tropicale tutta sua.

Siamo tutti reduci da una giornata di viaggio e l'intenzione non è di fare tardi. Cerchiamo di evitare i turisti occidentali che man mano diventano sempre più molesti — le inglesi in particolare, coi loro faccioni cresciuti a fish and chips — e le pozze di piscio di cavallo che, essendo l'unico mezzo di trasporto sull'isola, diventano un effetto collaterale non trascurabile.

Un posto interessante che ci promettiamo di rivalutare nei giorni seguenti è uno street food all'aperto, pieno di bancarelle con pesce fresco da grigliare sul momento. Elisa avrà scattato un centinaio di foto ai pesci esposti e un migliaio solo alle aragoste — ammetto che alcune erano veramente enormi. L'unico a provare qualcosa di strano alla fine sono sempre io: a una bancarella di dolci acquisto una crepe verde con ripieno di gelatina di riso e una pallina di pasta dolce ricoperta di semi di sesamo. Si chiamano Jajan Pasar, prodotti tradizionali indonesiani.

Durante il tragitto di ritorno passiamo davanti al ristorante dove abbiamo cenato e incontriamo il nostro cameriere, che ci saluta in spagnolo e ci chiede dove stiamo andando. "A descansar."

Ci separiamo dai nostri amici dopo aver fatto il piano per l'indomani mattina, poi — evitando le vacche libere per strada — rientriamo nella villa e crolliamo presto addormentati.

\pagestyle{empty} % disable numbers


\twoImagesLargeMargins{images/day8/PXL_20240816_041534888.jpg}
                      {images/day8/PXL_20240816_042132945.MP.jpg}

%coppia
\sixLandscapeImages{images/day8/PXL_20240816_041454119.jpg}
                   {images/day8/PXL_20240816_041433151.jpg}
                   {images/day8/PXL_20240816_044953711.jpg}
                   {images/day8/PXL_20240816_045026379.MP.jpg}
                   {images/day8/PXL_20240816_045040599.jpg}
                   {images/day8/PXL_20240816_045118434.jpg}

\twoImagesLargeMargins{images/day8/PXL_20240816_045630869.MP.jpg}
                      {images/day8/PXL_20240816_052323282.MP.jpg}
%

%coppia
\landscapeLargeMargin{images/day8/PXL_20240816_083751183.jpg}
\landscapeLargeMargin{images/day8/PXL_20240816_083729344.jpg}
%

%coppia
\portraitFullscreen{images/day8/PXL_20240816_112934774.MP.jpg}
\landscapeThinMarginCentered{images/day8/PXL_20240816_121112924.MP.jpg}
%

%coppia
\landscapeThinMarginCentered{images/day8/PXL_20240816_131943829.jpg}
\fourPortraitImages{images/day8/PXL_20240816_121516789.jpg}
                   {images/day8/PXL_20240816_121528574.jpg}
                   {images/day8/PXL_20240816_121536291.jpg}
                   {images/day8/PXL_20240816_120654036.MP.jpg}
%

%coppia 
\landscapeThinMarginCentered{images/day8/PXL_20240816_134909395.jpg}
\landscapeThinMarginCentered{images/day8/PXL_20240816_135458676.PORTRAIT.jpg}
%

%
\landscapeLargeMargin{images/day8/PXL_20240816_132309289.jpg}
%

\pagestyle{fancy} % enable numbers